\section{An opposing closed region for every elementary side}
\label{sec:area_law_elementary_side_opponents}

The actual plane cover supplies an opposing rectangle at the midpoint of
each elementary segment. Corner exclusion extends this contact to
containment of the whole segment; see
\cite[Section~11]{OpenAI2026AreaLaw}.

% Source: 10-geometry.tex, lines 154--181 and 299--310,
% prop:two-families, revision adc7f1241b42e322a6451854ab7e4b4c146bf78a.
% The conclusion is whole-segment containment, without uniqueness or labels.

\begin{theorem}[Existence of an opposing closed region]
    \label{thm:area_law_elementary_side_opponent}
    \lean{TNLean.PEPS.AreaLaw.Geometry.exists_elementarySide_opponent}
    \leanok
    \uses{def:area_law_dyadic_layers,def:area_law_fine_layer_indices,
        def:area_law_fine_belt_scales,def:area_law_actual_side_mask}
    Fix a common origin $o\in\mathbb R^2$, finite endpoint set $Z$,
    integer width $C_0\ge2$ and natural indices $k_0\le k$ with
    $k\ge50\,000\,000$. Let $z\in F_{k,\ell_k}$ be an actual
    fine-cell index, and let $e=[a_i,b_i]$ be any elementary
    segment of $Q_{k,z}=Q_{\ell_k}(z)$ under its actual midpoint
    subdivision relative to $k_0$. Then either
    \begin{align}
        e\subseteq\overline{N_{k_0}},
    \end{align}
    or there exist $h\ge k_0$ and $w\in F_{h,\ell_h}$ such that
    \begin{align}
        (k,z)\ne(h,w),\qquad
        e\subseteq\overline{Q_{\ell_h}(w)}.
    \end{align}
\end{theorem}
\begin{proof}\leanok
    \uses{thm:area_law_dummy_fine_cells_cover,thm:area_law_unique_fine_cell,
        thm:area_law_fine_cell_containment,thm:area_law_local_layer_indices,
        thm:area_law_fine_cells_disjoint,thm:area_law_dummy_layer_disjoint,
        thm:area_law_dummy_corner_elementary_side,
        thm:area_law_elementary_corner_endpoint,thm:area_law_dyadic_closures}
    The reference fine cell is nonempty, so its actual layer
    membership implies $Z\ne\varnothing$. Put $x=(a_i+b_i)/2$ and
    take the points in the ball of radius $10\,2^{\ell_k}$ about
    $x$ that lie outside the closed reference cell. Their closure
    contains $x$, since $x$ belongs to a side of that square.

    Every such point lies either in $N_{k_0}$ or in an actual
    fine cell of a layer $h\ge k_0$. Neighboring-layer locality
    restricts these latter indices to $k-1\le h\le k+1$.
    There are finitely many fine cells at these indices, and the
    reference cell is excluded by the choice of points. Taking
    the closure of this finite cover places $x$ either in the
    closed dummy neighborhood or in one distinct closed fine cell.
    In the dummy case, the finite union defining the neighborhood
    places $x$ in one contributing closed coarse cell.

    In either case the interiors of the selected rectangle and
    the reference square are disjoint. Since the midpoint of $e$
    belongs to the selected rectangle, its normal coordinate is
    a boundary coordinate of that rectangle. If the rectangle
    failed to contain either end of $e$, a rectangle corner
    would lie strictly inside $e$. The corresponding dummy-corner
    or actual fine-corner endpoint assertion excludes this.
    Thus the selected closed rectangle contains the whole segment.
    A contributing dummy rectangle belongs to $\overline{N_{k_0}}$,
    yielding the first alternative; a distinct actual fine cell
    yields the second.
\end{proof}

This assertion provides existence and whole-segment containment.
Opponent uniqueness, consistent primary or dummy labels, and the
isolated-star construction remain further obligations.
